% !TeX root = ../main.tex
\chapter{Grundlagen}
\label{cha:grundlagen}
% Nur, was zum Verständnis der Arbeit nötig ist. Empfehlung: dieses Kapitel
% ZULETZT zuschneiden — wer hier beginnt, investiert meist zu viel Zeit und
% Seiten. Technologien und Standards mit Primärquellen belegen (RFC,
% offizielle Dokumentation), siehe Beispiel unten.

\section{Beispielabschnitt}
\label{sec:grundlagen-beispiel}

Das \ac{http} ist ein zustandsloses Anwendungsprotokoll, dessen Semantik in
RFC~9110 festgelegt ist~\cite{rfc9110}. Für Python-spezifische Sprachkonstrukte
ist die offizielle Sprachreferenz die maßgebliche Quelle~\cite{python-reference}.
Als Beispiel eines Lehrbuchs dient~\cite{tanenbaum-betriebssysteme}.
